\documentclass[a4paper,12pt]{article}
\usepackage[utf8]{inputenc}
\usepackage[russian]{babel}
\usepackage{graphicx}
\usepackage{hyperref}
\usepackage{minted}
\usemintedstyle{friendly}
\usepackage{xcolor}
\usepackage{geometry}
\usepackage{amsmath}
\usepackage{enumitem}
\usepackage{float}
\usepackage{indentfirst}
\usepackage{caption}
\usepackage{amsfonts}
\usepackage{tikz}
\usetikzlibrary{shapes.geometric, arrows.meta, positioning}

\documentclass[14pt]{article}
\usepackage[utf8]{inputenc}
\usepackage[russian]{babel}
\usepackage{geometry}
\usepackage{minted}

\geometry{top=2cm, bottom=2cm, left=2.5cm, right=1.5cm}

% Настройки для листингов кода
\setminted{
    fontsize=\small,
    breaklines=true,
    autogobble=true,
    frame=single,
}

\begin{document}

\begin{titlepage}
    \centering
    % Верхний блок
    {\bfseries МИНИСТЕРСТВО НАУКИ И ВЫСШЕГО ОБРАЗОВАНИЯ РФ \\
    ФГБОУ ВО «Белгородский государственный технологический университет им. В.Г. Шухова»\par}
    
    \vfill
    
    % Центральный блок
    {\bfseries\Large Практическое задание №7\par}
    \vspace{0.5cm}
    {\large по дисциплине: Основы программирования и алгоритмизации \par}
    {\large тема: «Арифметика» \par}
    
    \vfill
    
    % Блок авторов
    \begin{flushright}
    \begin{tabular}{l}
        \large Выполнил: студент группы ВТ-252 \\
        \large Дронь Дмитрий Алексеевич \\
        \large Проверила: \\
        \large Стеценко Оксана Николаевна
    \end{tabular}
    \end{flushright}
    
    \vfill
    
    % Нижний блок
    {\large Белгород, 2026 г. \par}
\end{titlepage}

\newpage
\begin{center}
\section*{Цель работы}
\end{center}
Цель методических указаний организовать самостоятельную работу студента по выполнению индивидуальных заданий, направленных на закрепление навыков программирования на языке высокого уровня и способов решения математических задач с применением программы для работы с электронными таблицами MS Excel.
\newpage
\section*{Индивидуальное задание №7}
\textbf{Задача 5.} Для заданного получить таблицу первых простых чисел.
\newpage
\begin{center}
\section*{Блок-схема}
\end{center}
\begin{center}
\scalebox{0.75}{
\begin{tikzpicture}[
    >=Stealth,
    line/.style={draw,-Stealth,thick},
    block/.style={
        rectangle,
        rounded corners,
        draw,
        align=center,
        text width=3cm,
        minimum height=1cm
    },
    blockplain/.style={
        rectangle,
        draw,
        align=center,
        text width=3cm,
        minimum height=1cm
    },
    decision/.style={
        diamond,
        draw,
        align=center,
        aspect=2.2,
        text width=2.8cm,
        inner sep=2pt
    }
]

%------------------ Узлы ------------------
\node[block] (input) at (0,0)
{
is\_prime\\
x, flag = 1
};

\node[decision] (xcheck) at (0,-3)
{x < 2};

\node[blockplain] (flag0) at (0,-5.8)
{flag = 0};

\node[blockplain] (itwo) at (0,-8.3)
{i = 2};

\node[decision] (loopcheck) at (0,-11.3)
{i * i <= x \&\& flag == 1};

\node[decision] (modcheck) at (0,-14.5)
{x \% i == 0};

\node[blockplain] (flag0b) at (0,-17.3)
{flag = 0};

\node[blockplain] (iinc) at (0,-19.8)
{i++};

\node[block] (ret) at (0,-23)
{возврат flag};

%------------------ Стрелки ------------------
\draw (input) -- (xcheck);

% xcheck: да -> flag0
\draw
    (xcheck.south) -- ++(0,-0.2)
    node[pos=0.7,right]
    -- (flag0.north);


\draw (flag0.south) -- (itwo.north);

\draw (itwo.south) -- (loopcheck.north);

% loopcheck: -> modcheck
\draw
    (loopcheck.south) -- ++(0,-0.2)
    node[pos=0.7,right]
    -- (modcheck.north);

% loopcheck: -> ret (влево и вниз)
\draw[line]
    (loopcheck.west) -- ++(-3,0)
    -- ++(0,-10)
    -- ++(5.625,0)
    -- (ret.north);

% modcheck:а -> flag0b
\draw
    (modcheck.south) -- ++(0,-0.2)
    node[pos=0.7,right]
    -- (flag0b.north);



\draw (flag0b.south) -- (iinc.north);

% iinc -> loopcheck петля (вправо и вверх)
\draw[line]
    (iinc.east) -- ++(3.5,0)
    -- ++(0,10.5)
    -- ++(-5.15,0)
    -- (loopcheck.north);


\end{tikzpicture}
}
\end{center}

\begin{center}
\scalebox{0.80}{
\begin{tikzpicture}[
    >=Stealth,
    line/.style={draw,-Stealth,thick},
    block/.style={
        rectangle,
        rounded corners,
        draw,
        align=center,
        text width=3cm,
        minimum height=1cm
    },
    blockplain/.style={
        rectangle,
        draw,
        align=center,
        text width=3cm,
        minimum height=1cm
    },
    io/.style={
        trapezium,
        trapezium left angle=70,
        trapezium right angle=110,
        draw,
        align=center,
        text width=3cm,
        minimum height=1.1cm
    },
    decision/.style={
        diamond,
        draw,
        align=center,
        aspect=2.2,
        text width=3cm,
        inner sep=2pt
    }
]

%------------------ Узлы ------------------
\node[block] (input) at (0,0)
{
input\_t\\
*x, flag = 0,\\
tt
};

\node[io] (scantext) at (0,-3)
{Введите t};

\node[decision] (scanf1) at (0,-6)
{scanf("\%d", \&tt) != 1};

\node[blockplain] (flag0a) at (-4.5,-8)
{flag = 0};

\node[blockplain] (flag1a) at (4.5,-8)
{flag = 1};

\node[decision] (flagcheck1) at (0,-12)
{flag == 0};

\node[decision] (ttcheck) at (0,-15.5)
{tt <= 0 \textbar\textbar\ tt > MAXT};

\node[blockplain] (flag0b) at (-4.5,-17.5)
{flag = 0};

\node[blockplain] (flag2) at (4.5,-17.5)
{flag = 2};

\node[decision] (flagcheck2) at (0,-21.5)
{flag == 0};

\node[blockplain] (assignt) at (0,-24.3)
{*t = tt};

\node[block] (ret) at (0,-27)
{возврат flag};

%------------------ Стрелки ------------------
\draw (input) -- (scantext);
\draw (scantext) -- (scanf1);

% scanf1: нет -> flag0a (влево)
\draw
    (scanf1.west) -- ++(-1.87,0)
    node[midway,above]{нет}
    -- (flag0a.north);

% scanf1: да -> flag1a (вправо)
\draw
    (scanf1.east) -- ++(1.87,0)
    node[midway,above]{да}
    -- (flag1a.north);

% flag0a -> flagcheck1 (вниз и вправо)
\draw
    (flag0a.south) -- ++(0,-1.5)
    -- ++(4.5,0)
    -- (flagcheck1.north);

% flag1a -> flagcheck1 (вниз и влево)
\draw
    (flag1a.south) -- ++(0,-1.5)
    -- ++(-4.5,0)
    -- (flagcheck1.north);

% flagcheck1: да -> ttcheck
\draw[line]
    (flagcheck1.south) -- ++(0,-0.2)
    node[pos=0.7,right]{да}
    -- (ttcheck.north);


% ttcheck: нет -> flag0b (влево)
\draw
    (ttcheck.west) -- ++(-1.87,0)
    node[midway,above]{нет}
    -- (flag0b.north);

% ttcheck: да -> flag2 (вправо)
\draw
    (ttcheck.east) -- ++(1.87,0)
    node[midway,above]{да}
    -- (flag2.north);

% flag0b -> flagcheck2 (вниз и вправо)
\draw
    (flag0b.south) -- ++(0,-1.5)
    -- ++(4.5,0)
    -- (flagcheck2.north);

% flag2 -> flagcheck2 (вниз и влево)
\draw
    (flag2.south) -- ++(0,-1.5)
    -- ++(-4.5,0)
    -- (flagcheck2.north);

% flagcheck2: да -> assignt
\draw
    (flagcheck2.south) -- ++(0,-0.2)
    -- (assignt.north);


\draw (assignt.south) -- (ret.north);

\end{tikzpicture}
}
\end{center}
\newpage
\begin{center}
\scalebox{0.90}{
\begin{tikzpicture}[
    >=Stealth,
    line/.style={draw,-Stealth,thick},
    block/.style={
        rectangle,
        rounded corners,
        draw,
        align=center,
        text width=3.5cm,
        minimum height=1cm
    },
    blockplain/.style={
        rectangle,
        draw,
        align=center,
        text width=3.5cm,
        minimum height=1cm
    },
    decision/.style={
        diamond,
        draw,
        align=center,
        aspect=2.2,
        text width=2.6cm,
        inner sep=2pt
    }
]

%------------------ Узлы ------------------
\node[block] (input) at (0,0)
{
generate\_primes\\
t, primes[MAXT],\\
r.error\_flag = 0,\\
count = 0\\
x = 2
};

\node[decision] (countcheck) at (0,-4)
{count < t};

\node[decision] (primecheck) at (0,-7)
{is\_prime(x)};

\node[blockplain] (store) at (0,-10)
{
primes[count] = x;\\
count++;
};

\node[blockplain] (xinc) at (0,-12.8)
{x++};

\node[block] (ret) at (0,-16)
{возврат r};

%------------------ Стрелки ------------------
\draw (input) -- (countcheck);

% countcheck: -> primecheck
\draw
    (countcheck.south) -- ++(0,-0.2)
    -- (primecheck.north);

% countcheck: -> ret (влево и вниз)
\draw
    (countcheck.west) -- ++(-2.5,0)
    -- ++(0,-10.5)
    -- ++(4.325,0)
    -- (ret.north);

% primecheck: -> store
\draw
    (primecheck.south) -- ++(0,-0.2)
    -- (store.north);


% store -> xinc
\draw (store.south) -- (xinc.north);

% xinc -> countcheck петля (вправо и вверх)
\draw
    (xinc.east) -- ++(2,0)
    -- ++(0,10.5)
    -- ++(-3.9,0)
    -- (countcheck.north);

\end{tikzpicture}
}
\end{center}
\newpage
\section*{Программа}
\begin{minted}{c}
/* Летняя практика. Индивидуальное задание №7 «Арифметика»
   Выполнил студент группы ВТ-252 Дронь Д.А.
   Задание 5: Для заданного получить таблицу первых простых чисел.
*/
#include <stdio.h>

#define MAXT 10000

// Структура результата
typedef struct {
    int error_flag;
} Result;

// Проверка, является ли число простым
int is_prime(int x) {
    int flag = 1;

    if (x < 2) {
        flag = 0;
    }

    int i = 2;
    while (i * i <= x && flag == 1) {
        if (x % i == 0) {
            flag = 0;
        }
        i++;
    }

    return flag;
}

// Ввод t
int input_t(int *t) {
    int flag = 0;
    int tt;

    printf("Введите t: ");
    if (scanf("%d", &tt) != 1) {
        flag = 1;
    } else {flag = 0};

    if (flag == 0) {
        if (tt <= 0 || tt > MAXT) {
            flag = 2;
        } else {flag = 0};
    }

    if (flag == 0) {
        *t = tt;
    }

    return flag;
}

// Генерация первых t простых чисел
Result generate_primes(int t, int primes[MAXT]) {
    Result r;
    r.error_flag = 0;

    int count = 0;
    int x = 2;

    while (count < t) {
        if (is_prime(x)) {
            primes[count] = x;
            count++;
        }
        x++;
    }

    return r;
}

// Вывод таблицы простых чисел
void output_primes(int t, int primes[MAXT]) {
    int i = 0;
    printf("Первые %d простых чисел:\n", t);
    while (i < t) {
        printf("%d ", primes[i]);
        i++;
    }
    printf("\n");
}

// Основная функция
int main(void) {
    int t;
    int primes[MAXT];

    int err = input_t(&t);

    switch (err) {
        case 0:
            break;
        case 1:
            printf("Ошибка: некорректный формат ввода.\n");
            return 1;
        case 2:
            printf("Ошибка: t должно быть в диапазоне 1..%d.\n", MAXT);
            return 1;
        default:
            printf("Неизвестная ошибка ввода.\n");
            return 1;
    }

    Result r = generate_primes(t, primes);

    switch (r.error_flag) {
        case 0:
            output_primes(t, primes);
            break;
        default:
            printf("Неизвестная ошибка вычислений.\n");
            break;
    }

    return 0;
}

\end{minted}
\newpage
\section*{Тесты}
\begin{figure}[H]
    \centering
    \includegraphics[width=0.9\textwidth]{test.png}
    \caption{Результат выполнения программы}
\end{figure}
\newpage
\section*{Вывод}
В ходе выполнения работы мне удалось организовать самостоятельную работу по выполнению индивидуальных заданий, направленных на закрепление навыков программирования на языке высокого уровня и способов решения математических задач с применением программы для работы с электронными таблицами MS Excel.

\end{document}



